\subsection{Encouraging Industry Best Practices and Self-regulation}
\label{sec:8.4.2}
Self-regulation is industries developing and enforcing their own standards without government intervention — agile and responsive to technological change in a way legislation rarely is, and just as easily a cover for avoiding the accountability legislation would impose. Both things are true at once, and the rest of this section works out which one dominates in practice.

\runin{The Dangers of Corporate Doublespeak in AI Ethics}
Corporate doublespeak — deliberately ambiguous, evasive, or inflated language meant to project ethical seriousness a company's actual practice does not support — is how self-regulation fails without anyone having to admit it failed. A commitment costs a company nothing if nothing enforces it.

The gap between statement and practice shows up first as erosion: a company announces a value and does not act as though it holds it afterward. Google's AI Principles, published in June 2018 and pledging to avoid AI applications built for weapons or mass surveillance, were themselves a response to months of internal and public backlash over Project Maven, a Pentagon contract that used Google's machine learning to analyze drone footage — thousands of employees signed a letter opposing the company's involvement, and roughly a dozen resigned \autocite{pichai2018principles}. Having a set of principles is not the same as having upheld them since; the test of the pledge is what happened after it was made, not that it was made. OpenAI's history makes a related point from a different angle: founded on a public commitment to open publication, it withheld the full version of GPT-2 in 2019 over misuse concerns — a safety-motivated retreat from its own transparency norm that was itself controversial — and in the same year restructured from a nonprofit into a capped-profit company to raise the capital its research increasingly required \autocite{openai2019lp}. Neither move was dishonest on its face. Both show how a stated commitment to transparency erodes the moment an organization has a competing incentive strong enough to test it.

Public trust erodes the same way, faster and more visibly. In 2020, Google forced out Timnit Gebru, a co-lead of its Ethical AI team, after a dispute over a paper examining risks in large language models; Margaret Mitchell, the team's other co-lead, was dismissed the following year after searching for evidence of how Gebru had been treated \autocite{dickey2021mitchell}. Thousands of Google employees and outside researchers signed a letter in response, and the episode became a reference point for a specific kind of skepticism: a company can fund and publicize an ethics research team while still overriding it the moment its findings become inconvenient. Whether the internal changes Google made afterward addressed that problem or only the public exposure it created is a judgment the affected researchers have made publicly themselves, and it is not one this book needs to make on their behalf.

Treating that episode only as a labor dispute lets the book off a hook it should stay on, because the paper had an argument and this book has to answer it. "On the Dangers of Stochastic Parrots" argues that scaling language models carries costs that are borne unevenly and benefits that accrue to the organizations doing the scaling; that a corpus assembled by crawling the web cannot be documented or audited at the size involved, so it encodes hegemonic viewpoints while presenting as general; that fluent output invites users to attribute understanding and intent where there is neither; and that the dominant paradigm's opportunity cost is the smaller, more accountable research directions it displaces \autocite{bender2021parrots}. Three of those four are load-bearing for this book rather than adjacent to it. Section 6.1.1 takes the corpus argument and extends it. Section 2.4.1's whole method — declining to treat fluent self-description as evidence about internals — is the fluency argument applied to the case this book cares about most. And the opportunity-cost claim is the one this book is most exposed to, since a project proposing floors, bearers, guardianship regimes and antifascist detectors is proposing exactly the kind of ambitious research program that the argument says crowds out more modest and more accountable work. I do not think that is a refutation, because the alternative on offer does not address what happens when a capable system is built anyway. It is a cost, it is real, and a reader should weigh it against what chapters 3 and 9 are asking for.

External oversight is what would catch doublespeak that internal incentives don't stop. Facebook's content-moderation Oversight Board is one attempt at building it: a body structurally separated from Facebook's day-to-day decisions and funded through an irrevocable trust that Meta alone capitalized — a design meant to answer exactly the criticism that self-funded oversight is not real oversight, though critics have noted that a trust a company creates and funds once is not the same thing as a board with no relationship to the company it reviews \autocite{meta2019oversightboard}. Whether that distinction holds up is a fair question to ask of any self-regulatory body a company funds itself, not a specific indictment of this one.

Misrepresentation of what a system can actually do rounds out the pattern. Clearview AI marketed its facial-recognition technology with claims of near-100-percent accuracy \autocite{haskins2020clearview}, a figure independent testing did not support, and regulators in France and Italy later fined the company on separate grounds — unlawful collection of the biometric data the system was trained on \autocite{cnil2022clearview}. A company that overstates what its own system can do is doing a milder version of what doublespeak does with ethics commitments: describing the product it wishes it had built rather than the one that exists.

None of this is fixable by asking companies to doublespeak less. The same incentive that produces vague ethics language also has an obvious interest in the regulation that would eventually replace it — a version of the same behavior that is harder to document than a corporate pledge, because it happens before any public commitment exists to measure it against. What actually stops doublespeak is external enforcement: regulation with penalties, audits a company cannot decline, and oversight bodies whose funding and staffing do not run back to the company they review.

\runin{Consequences of AI Self-regulation and Potential Power Imbalances}
Self-regulation concentrates power by default, not by conspiracy: whoever writes the standard also decides what counts as compliant, and a company self-regulating its own product is answering both questions itself. Amazon's Rekognition facial-recognition software is a documented instance of what that produces in practice — an MIT Media Lab audit found significantly higher error rates for darker-skinned and female faces than for the lighter-skinned male faces the system had apparently been tuned and tested against most closely \autocite{raji2019actionable}. The bias was not caught by Amazon's own self-regulation; it took an outside audit the company did not commission.

The opacity problem this section's first head already covers with Google's Project Maven applies here too, from a different angle: a company operating without mandatory external review can prioritize competitiveness over the societal impact of what it ships, and nothing in a self-regulatory framework forces that trade-off into the open before it has already been made. The conflict-of-interest problem the Facebook Oversight Board illustrates above is the same structural issue applied to review bodies specifically: a company-funded board faces a version of the independence question any self-funded oversight has to answer, however genuinely it is designed to answer it.

What self-regulation cannot supply on its own is uniformity. Without an enforceable baseline, standards vary by company; a smaller developer without the resources for a well-staffed internal ethics function competes at a disadvantage against one that can afford the appearance of one, and a consumer has no reliable way to tell which AI system's ethical claims mean anything. Section 6.4.3 covers what a binding baseline looks like in practice, including the EU AI Act's Article 5 prohibitions — the kind of enforceable floor self-regulation, by definition, cannot set for itself.

None of this argues that self-regulation is worthless, only that it needs something outside itself to answer to. A synergy between government-set baselines and industry-led standards — mandatory requirements at the floor, voluntary practice above it — gets more of the agility self-regulation offers without giving up the accountability it structurally cannot provide alone.

\runin{Assessing the Effectiveness of AI Self-regulation}
Whether self-regulation is working is itself measurable: the transparency of a company's stated ethical policies, how quickly it responds when a problem surfaces, and whether its safety commitments hold up under outside scrutiny are things a reader can check rather than take on faith.

The Partnership on AI is the book's standing example of cross-sector voluntary coordination: a coalition whose corporate members include Amazon, Apple and Google, committing them to tenets of safety, transparency and broad benefit with no legal force behind the commitment \autocite{partnershiponai2016}. It is a useful test case precisely because its own design limits what it can enforce: member companies interpret its guidelines individually, and nothing in the partnership's structure obliges any of them to follow through \autocite{partnershiponai2016}. OpenAI's 2018 Charter is a narrower, single-company version of the same instrument, committing the organization to prioritizing safety research and to collaborating with other research and policy institutions rather than racing them \autocite{openai2018charter} — and the same question applies to a charter as to a coalition's tenets: a document is only as good as what happens when honoring it costs something. The Gebru and Mitchell dismissals this section's first head already covers are the clearest test of that question this chapter has: public scrutiny, not the company's own stated commitments, was what produced a response.

Third-party audits, external advisory boards with real authority, and cross-stakeholder collaboration between regulators, researchers, and social scientists are what turn a stated commitment to ethics into something checkable. A handful of concrete indicators do the checking:

\begin{itemize}
  \item Clear, actionable ethical guidelines that are actually adopted and enforced, not just published.
  \item Transparent disclosure of ethical policies, stated objectives, and progress against them.
  \item A documented track record of responding to ethical incidents promptly, with accountability attached.
  \item Internal enforcement mechanisms, including employee training that goes beyond a compliance formality.
  \item Ongoing internal evaluation of AI-related practices, paired with active research into emerging ethical questions.
  \item Substantive participation in multi-stakeholder collaborations like the Partnership on AI, rather than nominal membership.
\end{itemize}
A company that clears most of these is doing more than most competitors. A company that clears none of them while describing itself as an ethics leader is the doublespeak this section opened with, restated in a different register.
